% Appendix: detailed claim-incidence model structure.

% TODO: Model development, specifically precip. indices selection.
% Iterative cross validation, starting with the simplest precipitation indices,
% and gradually adding complexity based on predictive performance.
% Final decision was a balance of predictive performance, model complexity, and
% interpretability (important for the climate impact simulation use-case).

\begin{frame}{Total Damages Definition}
  \small
  \textit{Contract $i$, season-year $Q$ (e.g. Winter-2023):}
  \vspace{0.15cm}

  \begin{columns}[T,onlytextwidth]
    \begin{column}{0.31\textwidth}
      \centering
      $N_{iQ}$\\[+0.05cm]
      \footnotesize Number of claims
    \end{column}
    \begin{column}{0.31\textwidth}
      \centering
      $U_{iQ}^{k}$\\[0.05cm]
      \footnotesize Payout for claim $k$
    \end{column}
    \begin{column}{0.31\textwidth}
      \centering
      $A_i$\\[+0.05cm]
      \footnotesize Deductible
    \end{column}
  \end{columns}

  \vspace{0.25cm}

  Season-year $Q$ damage for contract $i$\,:
  \begin{equation*}
    S_{iQ} = \sum_{k=1}^{N_{iQ}} \left(A_i + U_{iQ}^{k}\right),
    \qquad \sum_{k=1}^{0} \left(A_i + U_{iQ}^{k}\right) = 0.
  \end{equation*}

  Total damage $T_{\mathcal{X},\mathcal{T}}$\,:
  contracts collection $\mathcal{X}$, year-seasons collection
  $\mathcal{T}$\,:
  \begin{equation*}
    T_{\mathcal{X},\mathcal{T}}
    = \sum_{Q \in \mathcal{T}} \sum_{i \in \mathcal{X}} S_{iQ}.
  \end{equation*}

  % \footnotesize
  % For example, $\mathcal{X}$ may contain all insured buildings in Vestland,
  % while $\mathcal{T}$ covers a 30-year period. We compare aggregate damage in
  % 1991--2020 with 2041--2070 under different CO$_2$ emission scenarios.
\end{frame}

\begin{frame}{Two precipitation indices survived validation}
  \begin{columns}[T,onlytextwidth]
    \begin{column}{0.53\textwidth}
      \begin{enumerate}
        \item \textbf{Absolute intensity}
          \begin{itemize}
            \item Mean precipitation across the three wettest days in
              a year-season
          \end{itemize}
        \item \textbf{Relative anomaly}
          \begin{itemize}
            \item Extreme-season rainfall relative to the local 95th
              percentile of daily precipitation
          \end{itemize}
      \end{enumerate}
    \end{column}
    \begin{column}{0.42\textwidth}
      \textbf{Interpretation}
      \begin{itemize}
        \item Damage requires sufficient water in absolute terms
        \item Local adaptation makes departures from normal
          conditions informative
      \end{itemize}

      Candidate indices were evaluated on unseen observations.
    \end{column}
  \end{columns}
  \sourceNote{Precipitation-index selection in the frequency model}
\end{frame}


\begin{frame}{Claim frequency model structure}
  %   \vspace{0.0cm}
  $N_{iQ}$: number of claims on contract $i$ in year-season $Q$;
  $l_{iQ}$: insured days that season.
  \centering
  \begin{align*}
    N_{iQ} \sim \ &\operatorname{NegBin}(\mu_{iQ},\, \phi), \\[0.3cm]
    \log\left(\frac{\mu_{iQ}}{l_{iQ}}\right)
    = \ &\beta_{\text{season}(Q)}, \\[0.10cm]
    \makebox[0.30\linewidth][l]{$\big[$\textit{\scriptsize property}$\big]$}
    +\ &X_i^{\text{prop.}}\beta^{\text{prop.}}
    + s(\text{build year}_i) + s(\text{tariff}_i), \\[0.1cm]
    \makebox[0.30\linewidth][l]{$\big[$\textit{\scriptsize topography}$\big]$}
    +\ &s(\text{TWI}_i)
    + s(\text{HAND}_i)
    + X_i^{\text{land use}}\beta^{\text{land use}}, \\[0.1cm]
    \makebox[0.30\linewidth][l]{$\big[$\textit{\scriptsize region}$\big]$}
    +\ &X_i^{\text{county}}\beta^{\text{county}}
    + X_i^{\text{munic.}}\beta^{\text{munic.}}_{\text{land use}}, \\[0.1cm]
    \makebox[0.30\linewidth][l]{$\big[$\textit{\scriptsize
    precipitation}$\big]$}
    +\ &s(\text{absolute precip}_{iQ})
    + s(\text{relative precip}_{iQ}).
  \end{align*}
  %   \vspace{0.3cm}
\end{frame}

\begin{frame}{Claim payment model structure}
  $U_{iQ}$: payout for a claim on contract $i$ in year-season $Q$.
  \centering
  \begin{align*}
    U_{iQ} \sim \ &\operatorname{Gamma}(\xi_{iQ},\, \nu), \\[0.3cm]
    \log(\xi_{iQ})
    = \ &\beta_{\text{season}(Q)}, \\[0.10cm]
    \makebox[0.30\linewidth][l]{$\big[$\textit{\scriptsize property}$\big]$}
    +\ &X_i^{\text{prop.}}\beta^{\text{prop.}}
    + s(\text{build year}_i)
    + \beta^{\text{prop.\ type}}_{\text{tariff}} \cdot \text{tariff}_i,
    \\[0.1cm]
    \makebox[0.30\linewidth][l]{$\big[$\textit{\scriptsize topography}$\big]$}
    +\ &s(\text{TWI}_i)
    + X_i^{\text{land use}}\, \beta^{\text{land use}}, \\[0.1cm]
    \makebox[0.30\linewidth][l]{$\big[$\textit{\scriptsize region}$\big]$}
    +\ &X_i^{\text{county}}\beta^{\text{county}}
    + X_i^{\text{munic.}}\beta^{\text{munic.}}_{\text{land use}},
    \\[0.1cm]
    \makebox[0.30\linewidth][l]{$\big[$\textit{\scriptsize
    precipitation}$\big]$}
    +\ &s(\text{absolute precip}_{iQ})
    + s(\text{relative precip}_{iQ}).
  \end{align*}
\end{frame}
